%=============================================================================!
% 17.4  FIRST-PRINCIPLES MOLECULAR DYNAMICS AND ITS COST                      !
%=============================================================================!
\section{First-principles molecular dynamics and its cost}
\label{sec:ip-aimd}

A map drawn from a survey is only as good as the survey, and a survey of
every field is too costly to repeat for each journey; the map is drawn
once and then used. A learned potential is such a map, and the survey is
electronic-structure theory, which gives the energy and forces of atoms
from their electrons (\cref{sec:pe-origin}). Molecular dynamics with
those forces at every step, first-principles or \emph{ab initio}
molecular dynamics, needs no fitted potential, and it is what a learned
potential is trained to imitate. This section measures what it costs,
with FHI-aims\cite{Blum2009}, which describes each electron's wave as a
sum of functions centred on the atoms, stored as tables of numbers, and
can run molecular dynamics itself.

\subsection*{One frame of LiC$_6$}

The 28 atoms of a $2\times2\times1$ cell of LiC$_6$, four copies of the
seven-atom cell of \cref{sec:ip-lic6}, taken from a run of MACE-MP-0 at
\SI{300}{\kelvin}, were given to FHI-aims with the PBE functional and its
standard `light' settings, the cheapest of its standard choices of basis
functions and integration grids, on four cores of an Apple M1 Max, once
for each of the three grids below. The electrons are found by a
self-consistent iteration: the potential each electron feels depends on
the density of all of them, the number per unit volume at each point,
which is built from their waves, which are in turn set by the potential;
so each cycle builds a potential from the density and a density from the
potential, until neither changes. This frame needed 16 or 17 cycles,
\SI{337}{\second} to \SI{385}{\second} in all: 14 to \SI{16}{\second} for
each ordinary cycle and 57 to \SI{60}{\second} for each of the last two,
which also computed the forces. A step of molecular dynamics starts from
the density of the step before, extrapolated, and needs fewer cycles. If
it took half the time of this calculation from scratch, an estimate that
a short run would test, a step would cost about \SI{192}{\second}, and a
picosecond of steps of \SI{1}{\femto\second} about \SI{53}{\hour} on these
four cores; the two cycles with forces do not shrink, so this is the
cheaper end. MACE-MP-0 with D3 takes \SI{214}{\milli\second} a step of
molecular dynamics for the same cell on the same machine's GPU, with other
runs sharing the machine, about 900 times less (and
\SI{72}{\milli\second} a call alone on the faster machine of
\cref{fig:ip-learned}a). The gap widens with size, since the cost of the
electronic structure grows faster than the number of atoms
(\cref{sec:pe-origin}).

\subsection*{A setting to converge}

LiC$_6$ is a metal. In a crystal each electron's wave repeats from one
cell to the next up to a shift of its phase, and the code finds the
electrons separately for each of a grid of such shifts, the $k$-points,
and averages over them; a metal, whose filled energy levels end
abruptly, needs a fine grid. A grid of $4\times4\times10$ shifts along
the three cell vectors has 160 points, but a shift and its reverse give
the same energies, which leaves 84 to compute; grids of $2\times2\times5$
and $3\times3\times8$ leave 12 and 37. The forces on the carbon atoms of
this frame, up to \SI{3.1}{\electronvolt\per\angstrom}, change by up to
\SI{0.16}{\electronvolt\per\angstrom} between the grids of 12 and 84
points and by up to \SI{0.28}{\electronvolt\per\angstrom} between those
of 37 and 84: they have not settled by 37 points, nothing yet shows that
they have by 84, and a denser grid is needed before these frames can
serve as a reference. A reference
calculation is converged in its own settings before anything is compared
with it, as a molecular dynamics run is converged in its time step
and length.

\subsection*{What is left to check}

The lithium studies of the next two sections rest on MACE-MP-0, and no
calculation here has yet compared its forces, or its barrier for a
lithium hop, with PBE for this system. Their numbers are therefore the
model's predictions. The check, single points at a converged grid on
frames of the runs and along lithium's paths between sites, is the first
step of the protocol of \cref{sec:ip-protocol} for any learned potential,
and the planned continuation of Part~V returns to it when it trains a model of its own on these data.

\begin{keyidea}
First-principles molecular dynamics computes the electrons at every step
and needs no fitted potential, at a cost of minutes a step for a few
dozen atoms, on this estimate about nine hundred times a foundation
model's here. Its settings,
such as the grid of $k$-points of a metal, are converged before it serves
as a reference, and a learned potential is checked against it on the
system at hand before its predictions are trusted.
\end{keyidea}

\notebook{17}{read the cost of a first-principles step and of a learned one}

\begin{selfcheck}
A first-principles step of 28 atoms costs \SI{192}{\second}. If the cost
grew as the cube of the number of atoms, as it does once finding the
eigenvectors dominates (\cref{sec:pe-origin}), roughly how long would a
step of 112 atoms take?
\end{selfcheck}
